\problemname{Coherency}

\illustration{0.3}{warhammer.jpg}{
    The miniatures are typically hand-painted by the players.
}
% Picture by Maarten Sijm, figurines painted by Alicia van Meer.

It is the 25th of March, 40025 CE
in the world of \emph{Battle Axe Player Clash 40,000} (BAPC40K).  % This is a name, not a number, so no math mode and no $\,$ as thousands-separator.
This futuristic table-top miniatures wargame
is played with endearing figurines called \emph{models}, each of which is placed on a circular \emph{base}.
The models are placed on a $100\,\text{km} \times 100\,\text{km}$ gaming board.
A collection of such models forms a \emph{coherent unit} if between any pair of models there is an unbroken chain of models
that have a Euclidean distance of at most two inches\footnote{%
    Recall that an inch equals $25.4\,\text{mm}$.
} between the edges of their bases.
Moreover,
if the unit contains seven models or more,
each model must be within two inches of at least \emph{two} other models.
Given the positions of a collection of models with varying base diameters,
determine whether they form a single coherent unit.

One can prove that for any valid input for this problem,
if the diameters of the circular bases differ from the given diameter by at most $10^{-5}\,\text{mm}$,
the coherency of a unit of models does not change.

\begin{Input}
The input consists of:
\begin{itemize}
\item One line with an integer $n$ ($2\leq n \leq 2\cdot 10^5$), the number of models.
\item $n$ lines, each with three integers $x$, $y$, and $d$
($0 \leq x, y \leq 10^8$, \\
$d\in\{25, 28, 32, 40, 50, 65, 80, 90, 100, 130, 160\}$),
describing a model that has its center coordinates at $(x,y)$ and a base diameter of $d$, all given in millimeters.
\end{itemize}
Each model (including the base) fits on the gaming board.

It is guaranteed that no two models are overlapping, but the models can touch.
\end{Input}

\begin{Output}
  If the $n$~models form a single coherent unit, output ``\texttt{yes}''.
  Otherwise, output ``\texttt{no}''.
\end{Output}

\remainingsamples

\vfill

\begin{figure}[h]
    \centering
    \includegraphics[width=0.85\textwidth]{samples.pdf}
    \caption{
        Illustration of the samples.
        Samples~1 and~4 are coherent. Sample~2 is not coherent, because the two models are too far away.
        Sample~3 is not coherent, because not all models are within two inches of two other models.
    }
\end{figure}
